\chapter{Lezione 6 --- Trasformata z: propriet\`a, tabelle fondamentali e funzione di trasferimento dei sistemi discreti}

\section{Richiamo: trasformata z della successione campionata}

Riprendiamo il risultato della lezione precedente. Se un segnale continuo causale \(f(t)\)
viene campionato con periodo \(T\), allora il segnale campionato in forma impulsiva si pu\`o scrivere come
\[
\hat f(t)=\sum_{k=0}^{+\infty} f(kT)\,\delta(t-kT).
\]

Applicando la trasformata di Laplace si ottiene
\[
\mathcal{L}\{\hat f(t)\}
=
\int_{0^-}^{+\infty}
\left(
\sum_{k=0}^{+\infty} f(kT)\,\delta(t-kT)
\right)e^{-st}\,dt
=
\sum_{k=0}^{+\infty} f(kT)e^{-skT}.
\]

Introducendo la variabile
\[
z=e^{sT},
\qquad\text{quindi}\qquad
e^{-sT}=z^{-1},
\]
si arriva alla definizione di trasformata z monolatera:
\[
X(z)=\mathcal{Z}\{x(k)\}
=
\sum_{k=0}^{+\infty} x(k)\,z^{-k}.
\]

Nel nostro caso \(x(k)=f(kT)\), quindi pi\`u rigorosamente si dovrebbe scrivere
\[
\mathcal{Z}\{f(kT)\}
=
\sum_{k=0}^{+\infty} f(kT)\,z^{-k}.
\]

\subsection*{Osservazione importante}

Nella trasformata z la dipendenza esplicita da \(T\) non compare pi\`u nella formula,
ma \`e \emph{nascosta} nella variabile \(z=e^{sT}\).
Per questo due successioni formalmente uguali possono corrispondere a segnali continui
campionati con periodi diversi.

\section{Trasformata z monolatera e bilatera}

Per segnali causali o comunque definiti per \(k=0,1,2,\dots\), si usa la trasformata z \textbf{monolatera}:
\[
X(z)=\sum_{k=0}^{+\infty} x(k)\,z^{-k}.
\]

Per segnali definiti su tutti gli interi, si pu\`o introdurre la trasformata z \textbf{bilatera}:
\[
X(z)=\sum_{k=-\infty}^{+\infty} x(k)\,z^{-k}.
\]

Nel seguito lavoriamo prevalentemente con la forma monolatera, perch\'e stiamo studiando
sistemi causali e sequenze con transitorio.

\section{Propriet\`a della trasformata z}

La trasformata z monolatera eredita molte propriet\`a della trasformata di Laplace.

\subsection{Linearit\`a}

Per \(\alpha,\beta\in\mathbb{R}\) oppure \(\mathbb{C}\),
\[
\mathcal{Z}\{\alpha x(k)+\beta y(k)\}
=
\alpha X(z)+\beta Y(z).
\]

\subsection{Convoluzione discreta}

Se definiamo la convoluzione discreta di due successioni causali come
\[
(x\ast y)(n)=\sum_{k=0}^{n} x(k)\,y(n-k),
\]
allora vale
\[
\mathcal{Z}\{x\ast y\}=X(z)Y(z).
\]

\subsection{Traslazione temporale}

Supponiamo di avere la successione
\[
x(k)\ \longleftrightarrow\ X(z)
\]
e consideriamo la successione anticipata di un passo
\[
\tilde x(k)=x(k+1).
\]

Allora
\[
\mathcal{Z}\{x(k+1)\}=zX(z)-zx(0).
\]

Questa propriet\`a \`e l'analogo discreto della propriet\`a della derivata
nella trasformata di Laplace.

\subsection{Teorema del valore iniziale}

Se la trasformata z esiste, allora
\[
x(0)=\lim_{z\to\infty} X(z).
\]

\subsection{Teorema del valore finale}

Se esiste il limite \(\lim_{k\to+\infty}x(k)\) e se sono soddisfatte le usuali ipotesi
sui poli di \((z-1)X(z)\), allora
\[
\lim_{k\to+\infty}x(k)
=
\lim_{z\to 1}(z-1)X(z).
\]

\section{Coppie fondamentali tra trasformata z, segnale campionato e trasformata di Laplace}

Raccogliamo alcune coppie molto usate negli esercizi.
Nella tabella seguente:
\begin{itemize}
    \item la prima colonna contiene la trasformata z della sequenza campionata;
    \item la seconda colonna il corrispondente segnale continuo causale;
    \item la terza colonna la relativa trasformata di Laplace.
\end{itemize}

\[
\renewcommand{\arraystretch}{1.5}
\begin{array}{|c|c|c|}
\hline
X(z) & x(t) & X(s) \\
\hline
1 & \delta(t) & 1 \\
\hline
\dfrac{z}{z-1} & 1(t) & \dfrac{1}{s} \\
\hline
\dfrac{Tz}{(z-1)^2} & t\,1(t) & \dfrac{1}{s^2} \\
\hline
\dfrac{T^2 z(z+1)}{2(z-1)^3} & \dfrac{1}{2}t^2\,1(t) & \dfrac{1}{s^3} \\
\hline
\dfrac{z}{z-e^{-aT}} & e^{-at}1(t) & \dfrac{1}{s+a} \\
\hline
\dfrac{Tze^{-aT}}{(z-e^{-aT})^2} & t\,e^{-at}1(t) & \dfrac{1}{(s+a)^2} \\
\hline
\dfrac{z\sin(\omega T)}{z^2-2\cos(\omega T)z+1}
& \sin(\omega t)\,1(t)
& \dfrac{\omega}{s^2+\omega^2}
\\
\hline
\dfrac{z\bigl(z-\cos(\omega T)\bigr)}{z^2-2\cos(\omega T)z+1}
& \cos(\omega t)\,1(t)
& \dfrac{s}{s^2+\omega^2}
\\
\hline
\dfrac{ze^{-aT}\sin(\omega T)}
{z^2-2e^{-aT}\cos(\omega T)z+e^{-2aT}}
& e^{-at}\sin(\omega t)\,1(t)
& \dfrac{\omega}{(s+a)^2+\omega^2}
\\
\hline
\end{array}
\]

\subsection*{Osservazione sui poli complessi coniugati}

Per il segnale
\[
e^{-at}\sin(\omega t)\,1(t),
\]
i poli della trasformata z sono
\[
z_{1,2}=e^{-aT}\bigl(\cos(\omega T)\pm j\sin(\omega T)\bigr)
=
e^{-aT}e^{\pm j\omega T}.
\]

Quindi:
\begin{itemize}
    \item il modulo dei poli \`e \(e^{-aT}\);
    \item l'angolo dei poli \`e \(\pm \omega T\).
\end{itemize}

\section{Sistema in forma di stato discreta}

Consideriamo ora un sistema lineare tempo-invariante nel discreto, in forma di stato:
\[
\begin{cases}
x(k+1)=Ax(k)+Bu(k),\\[4pt]
y(k)=Cx(k),
\end{cases}
\]
con condizione iniziale \(x(0)\).

Applicando la trasformata z monolatera all'equazione di stato, e usando la propriet\`a
di traslazione, si ottiene
\[
zX(z)-zx(0)=AX(z)+BU(z).
\]

Dall'equazione di uscita segue invece
\[
Y(z)=CX(z).
\]

Riordinando l'equazione di stato:
\[
(zI-A)X(z)=BU(z)+zx(0).
\]

Quindi
\[
X(z)=(zI-A)^{-1}BU(z)+(zI-A)^{-1}zx(0).
\]

Sostituendo nell'equazione di uscita:
\[
Y(z)=C(zI-A)^{-1}BU(z)+C(zI-A)^{-1}zx(0).
\]

\subsection{Interpretazione dei due termini}

I due termini hanno un significato fisico ben preciso:
\[
Y(z)=
\underbrace{C(zI-A)^{-1}BU(z)}_{\text{risposta forzata}}
+
\underbrace{C(zI-A)^{-1}zx(0)}_{\text{risposta libera}}.
\]

Il primo dipende dall'ingresso, il secondo dalla condizione iniziale.

\section{Matrice di trasferimento}

In assenza di condizioni iniziali, cio\`e ponendo \(x(0)=0\), resta
\[
Y(z)=C(zI-A)^{-1}BU(z).
\]

Da qui si definisce la \textbf{matrice di trasferimento}
\[
G(z)=C(zI-A)^{-1}B.
\]

Se il sistema ha anche un termine diretto \(D\), la formula generale diventa
\[
G(z)=C(zI-A)^{-1}B+D.
\]

Nel caso trattato negli appunti si assume \(D=0\).

\section{Calcolo di \((zI-A)^{-1}\) mediante aggiunta e determinante}

Per calcolare esplicitamente la matrice di trasferimento, pu\`o essere utile richiamare
la formula generale dell'inversa di una matrice invertibile.

Se \(M\in\mathbb{R}^{n\times n}\) \`e invertibile, allora
\[
M^{-1}=\frac{1}{\det(M)}\operatorname{Adj}(M),
\]
dove \(\operatorname{Adj}(M)\) \`e la matrice aggiunta di \(M\).

L'elemento \((i,j)\) dell'aggiunta si ottiene a partire dal complemento algebrico:
\[
c_{ij}=(-1)^{i+j}\det(M^{\ast}_{ji}),
\]
dove \(M^{\ast}_{ji}\) \`e il minore ottenuto eliminando la riga \(j\) e la colonna \(i\).

Applicando questa formula a \(M=zI-A\), si ottiene
\[
(zI-A)^{-1}
=
\frac{1}{\det(zI-A)}\operatorname{Adj}(zI-A).
\]

Quindi la matrice di trasferimento pu\`o essere scritta come
\[
G(z)=C(zI-A)^{-1}B
=
\frac{1}{\det(zI-A)}\,C\,\operatorname{Adj}(zI-A)\,B.
\]

\section{Caso SISO}

Nel caso \textbf{SISO} (Single Input Single Output),
\[
u(k)\in\mathbb{R},\qquad y(k)\in\mathbb{R},\qquad x(k)\in\mathbb{R}^n,
\]
e si pu\`o scrivere
\[
C=c^T,\qquad c\in\mathbb{R}^n,
\qquad
B=b,\qquad b\in\mathbb{R}^n.
\]

La funzione di trasferimento diventa quindi
\[
G(z)=c^T(zI-A)^{-1}b.
\]

Poich\'e compare l'inversa di una matrice polinomiale in \(z\),
la funzione di trasferimento \`e una \textbf{funzione razionale}:
\[
G(z)=
\frac{b_m z^m+b_{m-1}z^{m-1}+\dots+b_0}
{z^n+a_{n-1}z^{n-1}+\dots+a_1 z+a_0}.
\]

\subsection*{Osservazione}

Nel caso SISO, il denominatore coincide con il polinomio caratteristico
\[
\det(zI-A),
\]
mentre il numeratore dipende da \(A,B,C\).

\section{Equazione alle differenze associata alla funzione di trasferimento}

Sotto ipotesi di condizioni iniziali nulle, dalla relazione
\[
Y(z)=G(z)U(z)
\]
si ha
\[
\bigl(z^n+a_{n-1}z^{n-1}+\dots+a_0\bigr)Y(z)
=
\bigl(b_m z^m+b_{m-1}z^{m-1}+\dots+b_0\bigr)U(z).
\]

Formalmente, questa relazione porta a una equazione alle differenze.
Negli appunti compare una scrittura del tipo
\[
y(k+n)+a_{n-1}y(k+n-1)+\dots+a_0y(k)
=
b_m u(k+m)+b_{m-1}u(k+m-1)+\dots+b_0u(k).
\]

\subsection*{Forma causale pi\`u usata}

Per implementare il sistema come filtro digitale causale, \`e per\`o pi\`u comodo
scrivere la funzione di trasferimento in potenze di \(z^{-1}\):
\[
G(z)=
\frac{\beta_0+\beta_1 z^{-1}+\dots+\beta_m z^{-m}}
{1+\alpha_1 z^{-1}+\dots+\alpha_n z^{-n}}.
\]

In questo caso si ottiene direttamente la forma causale
\[
y(k)+\alpha_1 y(k-1)+\dots+\alpha_n y(k-n)
=
\beta_0 u(k)+\beta_1 u(k-1)+\dots+\beta_m u(k-m).
\]

Questa \`e l'equazione che definisce l'algoritmo di calcolo dell'uscita
quando l'ingresso attraversa un filtro con funzione di trasferimento \(G(z)\).

\section{Esempi di aggiunta e inversa di matrice}

Negli appunti compaiono due esempi espliciti di calcolo dell'aggiunta.

\subsection{Primo esempio: matrice \(2\times 2\)}

Sia
\[
M=
\begin{bmatrix}
3 & 6\\
7 & 9
\end{bmatrix}.
\]

Allora
\[
\det(M)=3\cdot 9-6\cdot 7=-15,
\]
e
\[
\operatorname{Adj}(M)=
\begin{bmatrix}
9 & -6\\
-7 & 3
\end{bmatrix}.
\]

Quindi
\[
M^{-1}
=
\frac{1}{-15}
\begin{bmatrix}
9 & -6\\
-7 & 3
\end{bmatrix}.
\]

\subsection{Secondo esempio: matrice \(3\times 3\)}

Sia
\[
M=
\begin{bmatrix}
5 & 2 & 0\\
6 & 6 & 1\\
0 & 3 & 2
\end{bmatrix}.
\]

Il determinante vale
\[
\det(M)=21,
\]
mentre l'aggiunta \`e
\[
\operatorname{Adj}(M)=
\begin{bmatrix}
9 & -4 & 2\\
-12 & 10 & -5\\
18 & -15 & 18
\end{bmatrix}.
\]

Quindi
\[
M^{-1}
=
\frac{1}{21}
\begin{bmatrix}
9 & -4 & 2\\
-12 & 10 & -5\\
18 & -15 & 18
\end{bmatrix}.
\]

\section{Sintesi della lezione}

I punti chiave della lezione sono:
\begin{enumerate}
    \item la trasformata z di una successione causale si definisce come
    \[
    X(z)=\sum_{k=0}^{+\infty}x(k)z^{-k};
    \]
    \item la dipendenza dal periodo di campionamento \(T\) scompare formalmente,
    ma resta implicita nella relazione \(z=e^{sT}\);
    \item la trasformata z monolatera ha propriet\`a di linearit\`a, convoluzione,
    traslazione nel tempo, valore iniziale e valore finale;
    \item esistono tavole di corrispondenza tra segnali continui causali,
    trasformate di Laplace e trasformate z delle rispettive sequenze campionate;
    \item per un sistema discreto in forma di stato
    \[
    \begin{cases}
    x(k+1)=Ax(k)+Bu(k),\\
    y(k)=Cx(k),
    \end{cases}
    \]
    si ottiene
    \[
    Y(z)=C(zI-A)^{-1}BU(z)+C(zI-A)^{-1}zx(0);
    \]
    \item la matrice di trasferimento \`e
    \[
    G(z)=C(zI-A)^{-1}B;
    \]
    \item nel caso SISO, \(G(z)\) \`e una funzione razionale;
    \item dalla funzione di trasferimento si ricava l'equazione alle differenze
    che implementa il filtro discreto;
    \item il calcolo esplicito di \(G(z)\) pu\`o essere fatto usando
    \[
    (zI-A)^{-1}=\frac{1}{\det(zI-A)}\operatorname{Adj}(zI-A).
    \]
\end{enumerate}